\documentclass[a4paper,12pt]{article}

\usepackage[T2A]{fontenc}
\usepackage[utf8]{inputenc}
\usepackage[russian]{babel}
\usepackage{amsmath,amssymb,mathtools}
\usepackage{PTSans}
\renewcommand{\familydefault}{\sfdefault}
\usepackage{icomma}
\usepackage[a4paper,margin=20mm,headheight=25pt]{geometry}
\usepackage{microtype}
\usepackage{xcolor}
\usepackage{tikz}
\usetikzlibrary{calc,arrows.meta,positioning}
\usepackage{tkz-euclide}
\usepackage{booktabs,tabularx}
\usepackage{enumitem}
\usepackage{parskip}
\usepackage{fancyhdr}
\usepackage{setspace}
\usepackage[hidelinks]{hyperref}

\definecolor{accent}{HTML}{2A7775}
\definecolor{darkaccent}{HTML}{1B5553}
\definecolor{textgray}{HTML}{303536}
\definecolor{lightgray}{HTML}{6C7675}
\color{textgray}
\onehalfspacing
\setlength{\parindent}{0pt}
\setlength{\parskip}{0.5em}
\setlist[itemize]{leftmargin=1.45em,topsep=0.2em,itemsep=0.15em}
\setlist[enumerate]{leftmargin=1.65em,topsep=0.35em,itemsep=0.5em}
\pagestyle{fancy}
\fancyhf{}
\fancyhead[L]{\small\color{darkaccent}Геометрия: углы и отрезки}
\fancyhead[R]{\small\color{darkaccent}Кирилл Зиновьев}
\fancyfoot[C]{\small\color{lightgray}\thepage}
\renewcommand{\headrulewidth}{0.4pt}
\renewcommand{\headrule}{\hbox to\headwidth{\color{accent}\leaders\hrule height\headrulewidth\hfill}}
\newcommand{\headingrule}{\vspace{-0.45em}{\color{accent}\hrule height 0.55pt}\vspace{0.55em}}
\newcommand{\key}[1]{\textbf{\color{darkaccent}#1}}
\newcommand{\task}[1]{\item \textbf{#1}}

\begin{document}
% Титульный лист
\thispagestyle{empty}
\begin{center}
\vspace*{22mm}
{\Large\color{darkaccent}\bfseries ЧЕК-ЛИСТ}\\[9mm]
{\huge\bfseries Углы и отрезки}\\[3mm]
{\LARGE\bfseries Биссектриса, смежные и вертикальные углы}\\[10mm]
{\large Геометрия, 7 класс}\\[3mm]
{\large 10 октября 2026 года}\\[17mm]
{\large\bfseries Кирилл Зиновьев}\\[12mm]
\begin{minipage}{0.79\textwidth}
\centering\normalsize
Как распознавать отношения между углами, использовать их свойства
и находить длины отрезков по чертежу и условию задачи.
\end{minipage}
\vfill
{\large Лёвин Артём Александрович}\\[2mm]
{\normalsize эксклюзивно для Кирилла Зиновьева}\\[11mm]
\begin{tikzpicture}
\draw[accent,line width=1.2pt]
(-6.1,0) .. controls (-3.8,0.75) and (-1.9,-0.55) .. (0,0)
.. controls (2.2,0.75) and (4.15,-0.45) .. (6.1,0.2);
\end{tikzpicture}
\end{center}
\clearpage

\section*{1. Обозначение угла и биссектриса}
\headingrule
\key{Обозначение угла.} В записи $\angle DCL$ буква $C$ в середине означает \emph{вершину} угла. Угол измеряют в градусах; в числовом ответе знак $^\circ$ обязателен.

\key{Биссектриса угла} --- луч, выходящий из его вершины и делящий угол на два равных угла.

\begin{center}
\begin{tikzpicture}[scale=1]
\tkzDefPoint(0,0){C}
\tkzDefPoint(4.2,0){L}
\tkzDefPoint(-2.351,3.236){D}
\tkzDefPoint(1.044,2.049){M}
\tkzDrawSegments[color=textgray,line width=1pt](C,L C,D)
\tkzDrawSegments[color=accent,line width=1.5pt](C,M)
\tkzDrawPoints[size=2.5,fill=darkaccent](C,D,L,M)
\tkzLabelPoints[below](C,L)
\tkzLabelPoints[above](D,M)
\node at (1.9,0.43) {\small $63^\circ$};
\node at (-0.06,1.32) {\small $63^\circ$};
\end{tikzpicture}
\end{center}
\textbf{Пример.} $\angle DCL=126^\circ$, луч $CM$ --- биссектриса.
Тогда
\[
\angle DCM=\angle MCL=\frac{126^\circ}{2}=63^\circ.
\]
\key{Проверьте себя:} биссектриса делит величину угла \emph{пополам}; обе полученные части должны быть равны.

\section*{2. Смежные углы}
\headingrule
\key{Смежные углы} имеют общую сторону, а две другие стороны являются противоположными лучами. \key{Их сумма равна $180^\circ$.}
\[
\angle AOB+\angle BOC=180^\circ,
\qquad \text{если лучи }OA\text{ и }OC\text{ противоположны.}
\]
\textbf{Пример.} Один из смежных углов в $3$ раза больше другого.
Обозначим меньший угол через $x$, больший --- через $3x$:
\[
x+3x=180^\circ,\qquad 4x=180^\circ,\qquad x=45^\circ.
\]
Следовательно, углы равны $45^\circ$ и $135^\circ$.

\key{Важный вывод.} Два смежных угла не могут быть тупыми: каждый тупой угол больше $90^\circ$, а сумма двух таких углов превысила бы $180^\circ$.

\clearpage
\section*{3. Вертикальные углы}
\headingrule
При пересечении двух прямых образуются четыре угла.
\key{Вертикальные углы} расположены друг напротив друга; стороны одного угла являются противоположными лучами к сторонам другого. \key{Вертикальные углы равны.}

\begin{center}
\begin{tikzpicture}[scale=1.05]
\tkzDefPoint(-2.7,-1.5){A}
\tkzDefPoint(2.7,1.5){C}
\tkzDefPoint(-2.3,1.65){B}
\tkzDefPoint(2.3,-1.65){D}
\tkzDefPoint(0,0){O}
\tkzDrawSegments[color=textgray,line width=1pt](A,C B,D)
\tkzDrawPoints[size=2.2,fill=darkaccent](O)
\tkzLabelPoints[left](A,B)
\tkzLabelPoints[right](C,D)
\tkzLabelPoints[above right](O)
\node at (-1.05,0) {\small $\alpha$};
\node at (1.07,0) {\small $\alpha$};
\node at (0,0.9) {\small $\beta$};
\node at (0,-0.9) {\small $\beta$};
\end{tikzpicture}
\end{center}
Для прямых $AC$ и $BD$:
\[
\angle AOB=\angle COD,\qquad
\angle BOC=\angle DOA.
\]
Каждая пара соседних смежных углов в сумме даёт $180^\circ$.

\textbf{Пример.} Сумма двух вертикальных углов равна $86^\circ$.
Каждый из них равен $86^\circ:2=43^\circ$.
Смежный с ним угол составляет
\[
180^\circ-43^\circ=137^\circ.
\]

\section*{4. Как выбрать нужное свойство}
\headingrule
\begin{center}
\renewcommand{\arraystretch}{1.28}
\begin{tabularx}{\linewidth}{@{}>{\bfseries}p{36mm}X X@{}}
\toprule
Углы & Как узнать & Что использовать\\
\midrule
Смежные & Есть общая сторона; две другие --- противоположные лучи. & Сумма $180^\circ$.\\
Вертикальные & Стороны попарно противоположны; углы друг напротив друга. & Углы равны.\\
Части угла с биссектрисой & Луч из вершины делит угол на две равные части. & Каждая часть --- половина исходного угла.\\
\bottomrule
\end{tabularx}
\end{center}
\key{Порядок решения:} аккуратный чертёж $\rightarrow$ подписи вершин
$\rightarrow$ определение свойства $\rightarrow$ вычисления $\rightarrow$ ответ со знаком $^\circ$.

\clearpage
\section*{5. Отрезки на одном луче}
\headingrule
Если точки $A$ и $B$ лежат на одном луче с началом $O$, то более удалённая от $O$ точка соответствует большему расстоянию.
\[
AB=\bigl|OA-OB\bigr|.
\]
\textbf{Пример.} $OA=10{,}3$ см, $OB=2{,}4$ см. Следовательно, точка $B$ ближе к $O$, чем точка $A$:
\[
AB=10{,}3-2{,}4=7{,}9\ \text{см}.
\]
\key{Проверка чертежа:} длины на рисунке должны качественно соответствовать условию. Если $OA>OB$, изображайте $A$ дальше от $O$, чем $B$.

\section*{6. Деление отрезка в данном отношении}
\headingrule
Запись $AM:MB=3:7$ означает: весь отрезок $AB$ состоит из $3+7=10$ равных долей, причём $AM$ содержит $3$ доли, а $MB$ --- $7$ долей.

\textbf{Пример.} На отрезке $AB$ расположены точки $M$ и $N$, причём
\[
AM:MB=3:7,\qquad AN:NB=7:3,\qquad MN=38\ \text{см}.
\]
\begin{center}
\begin{tikzpicture}[x=1cm,y=1cm]
\tkzDefPoint(0,0){A}
\tkzDefPoint(3,0){M}
\tkzDefPoint(7,0){N}
\tkzDefPoint(10,0){B}
\tkzDrawSegments[color=textgray,line width=1.1pt](A,B)
\tkzDrawPoints[size=3,fill=darkaccent](A,M,N,B)
\tkzLabelPoints[below](A,M,N,B)
\draw[accent,thin] (0,0.25)--(0,0.42) (3,0.25)--(3,0.42) (7,0.25)--(7,0.42) (10,0.25)--(10,0.42);
\node at (1.5,0.52) {\small 3 доли};
\node at (5,0.52) {\small 4 доли};
\node at (8.5,0.52) {\small 3 доли};
\end{tikzpicture}
\end{center}
По первому отношению $AM=\frac{3}{10}AB$; по второму $AN=\frac{7}{10}AB$.
Поэтому
\[
MN=AN-AM=\frac{7}{10}AB-\frac{3}{10}AB=\frac{4}{10}AB.
\]
Отсюда $\frac{4}{10}AB=38$, значит $AB=38:\frac{4}{10}=95$ см.

\key{Важная тонкость:} для двух точек условия деления задают \emph{отдельно}; по одной фразе «точки делят отрезок в отношении $3:7$» нельзя однозначно установить их положение.

\section*{7. Перед сдачей работы}
\headingrule
\begin{itemize}
\item[$\square$] На чертеже указаны вершины углов и точки отрезка.
\item[$\square$] Острый угол выглядит острым, тупой --- тупым; большие отрезки не показаны короче меньших.
\item[$\square$] Для смежных углов использована сумма $180^\circ$, для вертикальных --- равенство.
\item[$\square$] Вычисления записаны даже тогда, когда ответ кажется очевидным.
\item[$\square$] В ответах стоят знаки градусов или единицы длины.
\end{itemize}

\clearpage
\section*{Тренировочный блок. Путь А}
\headingrule
\textit{10 заданий по возрастанию сложности. Делайте чертёж там, где он помогает решению. Для каждой задачи запишите вычисления и итоговый ответ.}

\begin{enumerate}[label=\textbf{\arabic*.},itemsep=0.95em]
\item Угол $DCL$ равен $126^\circ$. Луч $CM$ --- его биссектриса. Найдите $\angle MCL$.

\item $\angle KLM=56^\circ$. Луч $LN$ делит этот угол пополам. Найдите $\angle KLN$.

\item Углы $AOB$ и $BOC$ смежные. Известно, что $\angle AOB=105^\circ$. Найдите $\angle BOC$.

\item Один из смежных углов в $5$ раз больше другого. Найдите оба угла.

\item Один из смежных углов на $38^\circ$ больше другого. Найдите оба угла.

\item Прямые $AC$ и $BD$ пересекаются в точке $O$. Известно, что $\angle AOB=72^\circ$. Найдите $\angle COD$, $\angle BOC$ и $\angle DOA$.

\item Прямые $KS$ и $LM$ пересекаются в точке $O$. Сумма вертикальных углов $KOL$ и $SOM$ равна $146^\circ$. Найдите $\angle LOS$.

\item На одном луче с началом в точке $O$ лежат точки $A$ и $B$. Известно, что $OA=12{,}7$ см, а $OB=4{,}9$ см. Какая из точек ближе к $O$? Найдите $AB$.

\item Углы $AOB$ и $BOC$ смежные, причём $\angle AOB=2\angle BOC$. Луч $OD$ --- биссектриса угла $AOB$. Найдите $\angle DOB$.

\item Точки $M$ и $N$ лежат внутри отрезка $AB$. Известно, что $AM:MB=3:7$, $AN:NB=7:3$ и $MN=28$ см. Найдите длину $AB$.
\end{enumerate}

\clearpage
\section*{Ответы к тренировочному блоку}
\headingrule
\begin{center}
\renewcommand{\arraystretch}{1.55}
\begin{tabularx}{0.92\linewidth}{@{}c X@{}}
\toprule
\textbf{№} & \textbf{Ответ}\\
\midrule
1 & $63^\circ$\\
2 & $28^\circ$\\
3 & $75^\circ$\\
4 & $30^\circ$ и $150^\circ$\\
5 & $71^\circ$ и $109^\circ$\\
6 & $\angle COD=72^\circ$; $\angle BOC=108^\circ$; $\angle DOA=108^\circ$\\
7 & $107^\circ$\\
8 & Ближе точка $B$; $AB=7{,}8$ см\\
9 & $60^\circ$\\
10 & $70$ см\\
\bottomrule
\end{tabularx}
\end{center}
\vspace{1.2em}
\textit{Самопроверка.} Для задач об углах проверьте, что смежные значения дают сумму $180^\circ$, а вертикальные совпадают. Для задач об отрезках проверьте порядок точек на рисунке и единицы измерения.

\end{document}
